\section{A common denominator setup and its numerical table}

% BEGIN F78D1
\paragraph{F78D1. The degree-one prime in the family setup.}
The Chebotarev hypothesis in SSC17 holds in the convention of the
original family. For an open normal subgroup $U$ of the absolute
Galois group of a number field $E$, let $F/E$ be its fixed field.
It is finite Galois, and its fixing subgroup is exactly $U$.
Given $g\in G_E$ and finitely many excluded primes $A$ of $E$,
use the construction of CAP6 in a finite Galois extension
$F'/\mathbb Q$ containing $F$. Exclude the rational primes below
$A$ and those ramifying in $F'$. A prime in the Frobenius class of
$g|_{F'}$ can be conjugated so that its Frobenius equals that
element. The prime $v$ below it in $E$ has degree one, since
$g$ fixes $E$ and the $p$-power map then fixes the residue field.
It also has ramification index one over $\mathbb Q$. For the
extension $B$ of this prime to the chosen algebraic closure, the
inertia over $E$ fixes $F$ because $F/E$ is unramified there.
CAP6 gives an arithmetic Frobenius lift $\sigma$ with the
specified residue power map and with $\sigma|_F=g|_F$. Thus
$I_B\subset U$ and $\sigma g^{-1}\in U$, with $v\notin A$.
These include the inertia and ramification clauses needed by
the canonical hub construction and by the six Selmer conditions.
% END F78D1

% BEGIN F78D2
\paragraph{F78D2. The full original family setup with its uniform cutoff.}
Assume the original $\mathrm{Hypotheses}_2$, a family $\mathrm{Fm}$
with exponent $\varepsilon$, real numbers $C,C_2$, and a natural
number $C_T$. Suppose $\mathrm{TwGood}(\mathrm{Fm})$ and
$\mathrm{FactsF}(\mathrm{Fm},C,C_T)$. There exists $x_0$ such
that, for every $x\ge x_0$ with $x>\mathrm{Fm}.x_1$, every
$i\in\mathrm{Fm}.\mathrm{Ix}(x)$ satisfying the stated odd
Tate--Shafarevich finiteness premise on $D=\mathrm{Fm}.D(x,i)$,
and every finite exceptional set $S$ of natural numbers, one has
a Wach system $H$, its functor $F$, its comparison, and a hub
$X$ on that $D$ satisfying
\[
 X.\mathrm{exc}=S,\qquad \mathrm{AdmJ}(L),\qquad
 \mathrm{Setup}(L,\mathrm{nch},H,X,C_2).
\]
Here $L$ and $\mathrm{nch}$ are the size and character-choice data
of this very family member.

Take the maximum of the cutoff for admissible parameters at the
eighth family parameter (index $7$, default zero), the small-prime
cutoff at constant one, and the cutoff supplied by $\mathrm{TwGood}$.
The family has a member at zero, whose positive scale exponent equals
$\varepsilon$, so the parameter estimate applies. The family
identifies the local parameter with this indexed bound and identifies
the input scale with $x$. At the chosen member, apply L78W13 to the
Bhatt--Scholze field of $\mathrm{Hypotheses}_2$ on $D$. This
produces the three compatible local records on $D$.
Supply them and F78D1 to SSC17. More explicitly, the good twist
gives the index for DFD9, where the exceptional-set function is
constant $S$ and the error function is zero. The resulting hub
uses that index, has the specified exceptional set, and has the
canonical point family for $H$ and $\mathrm{rowC}(C_2)$.
SSC15--SSC16 and DFD10--DFD11 give killing on the representation
of this hub. Its remaining setup fields are the family facts,
canonical record, admissible parameters, polynomial parameter
bounds, $8\le L.\mathrm{sc}.L$, the splitting bound at least
$60$, and the supplied odd finiteness. The admissible parameters
also give $\mathrm{AdmJ}(L)$ by the member's facts. The maximum
was chosen before $x,i$ and $S$. This gives both source statements
\texttt{exists\_setup\_large} and \texttt{exists\_hub\_data\_family}
with their original quantifier order.
% END F78D2

% BEGIN F78D3
\paragraph{F78D3. All five numerical clauses on that same hub.}
In F78D2 take $C_2\ge1$. For any real $C_s$ with
$C_s\ge\max(3C_2+1,22365)$, the original five numerical table
clauses at constant $C_s$ hold on the very hub $X$ produced there,
at the same cutoff and for the same prescribed exceptional set.

Indeed the canonical record in the setup supplies
$\mathrm{DataAD}(L,H,X,R)$ for
$R=\mathrm{rowC}(C_2)=(C_2,320,25,23,32,32)$.
Its derived constants are
\[
 K_1=320+48=368,\qquad K_2=6\cdot320+12+6\cdot25=2082,
 \qquad K_3=3\cdot23\cdot320+12\cdot23+9=22365.
\]
The fixed inequalities in $R.\mathrm{Le}(C_s)$ now read
$K_1,K_2+1,K_3,6\cdot320+1,2(3\cdot23+1)\le C_s$,
in addition to $1\le C_s$, $3C_2+1\le C_s$, $1\le C_2$
and nonnegativity of the three fixed row constants.
The displayed bound on $C_s$ proves each. Apply the row A, B, C
and D estimates in the preceding section with this $R$ and with
the representation, column and lengths of $X$. For the fifth
clause, use the polynomial parameter bound $\mathrm{par67}$ from
this same setup. This yields precisely the field
$\mathrm{num}$ in the family conclusion. All five entries and
the killing predicate of F78D2 use the same $X$. The other bounds
on that representation still have to be supplied to construct the
full family conclusion.
% END F78D3
